\chapter{Conclusion générale}

Les parties finales du projet montrent une chaîne méthodologique complète allant de la construction des variables à l'interprétation du modèle final et à l'analyse opérationnelle des événements extrêmes.

Le modèle final retenu, ExtraTreesRegressor, conserve de bonnes performances sur l'ensemble de test, avec une MAE d'environ 0.467, un RMSE d'environ 0.986 et un R² d'environ 0.893. Ces résultats indiquent une bonne capacité de généralisation sur une période future indépendante.

L'interprétabilité confirme la cohérence physique du modèle. Les variables les plus importantes sont liées au CSI, à la variation récente de la production normalisée, aux variations de l'irradiance et au potentiel théorique d'irradiance directe sous ciel clair, en particulier \textit{clear\_sky\_bni}. Les analyses par feature importance, SHAP et permutation importance convergent donc vers la même conclusion : le modèle s'appuie principalement sur des signaux pertinents pour détecter les changements rapides de conditions solaires.

L'analyse des rampes critiques montre que le modèle ne se limite pas à de bonnes performances moyennes en régression. Avec le seuil Q90, les alertes sont plus fiables, avec une précision d'environ 81.9 \%. Avec le seuil optimisé, le modèle détecte davantage d'événements critiques, avec un recall d'environ 85.5 \%, au prix d'un plus grand nombre de fausses alertes.
